\section{The frame offset \texorpdfstring{$\beta_{\rm sym}$}{beta\_sym}}\label{sec:beta}

Most physical readouts of the Companion depend on one number, the \emph{frame offset} \(\beta_{\rm sym}\approx0.549\) between the matter
frame and the light frame.  This section states how it enters, why it is an input, and what the best current reading of it is
and is not.  Its derivation is the open problem OP1.

\subsection{Definition}

In the frame picture the matter frame and the tachyon frame sit at rapidities \(+\eta\) and \(-\eta\) on either side of the light frame
(Companion, theorem \texttt{thm:symmetric-frame}).  Write
\[
\beta_{\rm sym}=\tanh\eta,\qquad \gamma_L=\cosh2\eta,\qquad\text{so that}\qquad
\beta_{\rm sym}=\sqrt{\frac{\gamma_L-1}{\gamma_L+1}} .
\]
Thus \(\gamma_L\) is the Lorentz factor between the two opposite frames.  The Companion sets
\[
\gamma_L=\frac{\gamma_{\rm BI}}{\gamma_{\rm string}},\qquad
\gamma_{\rm string}=\frac{\ln2}{\pi\sqrt3}\approx0.12738,\qquad \gamma_{\rm BI}\approx0.2375,
\]
which gives \(\gamma_L\approx1.8644\), \(\beta_{\rm sym}\approx0.5493\) and \(\eta\approx0.6175\).  Both numbers on the right are values of the
Barbero--Immirzi parameter of loop quantum gravity, taken from the literature.  So \(\beta_{\rm sym}\) is \emph{put in}; nothing in
Sections~\ref{sec:kernel}--\ref{sec:generation} fixes it.

\subsection{Two candidate values that do not apply}

Two structural values had been proposed for \(\beta_{\rm sym}\): the ``Hopf value'' \(1/\sqrt3\approx0.577\), from half the angle of a
\(\mathbb Z_3\)-symmetric spinor on \(S^3\), and a ``Tribonacci'' condition \(\gamma_L=1/\beta_{\rm sym}\).  A frame audit (layer
v4.9.0-r61) showed that both compare quantities from different frames or of different kind.  The Hopf value describes a pure state read
from the light pole, which at \(\beta^2=\tfrac13\) has populations \((\tfrac23,\tfrac13)\), while the matter frame has \((\tfrac34,\tfrac14)\);
the same symbol labels two different states.  The Tribonacci condition equates a readout of the matter frame with the offset of the
tachyon in its own measure, i.e.\ \(\cosh2\eta=e^\eta\).  Neither is a frame-consistent fixed point.  OP1 therefore asks for the measured
value itself, not for a correction of these.

\subsection{The counting reading}

The two input values are solutions of one equation.  In the U(1) treatment of the black-hole horizon~\cite{abck,domagala,meissner}
the parameter \(\gamma\) solves
\[
\sum_j g(j)\,e^{-2\pi\gamma\sqrt{j(j+1)}}=1,
\]
with the sum over the allowed spins and \(g\) a degeneracy.

\begin{proposition}[Countings]\label{prop:countings}
Spin \(j=\tfrac12\) alone with \(g=2\) gives \(\gamma_{\min}=\ln2/(\pi\sqrt3)=\gamma_{\rm string}\).  All spins with \(g=2\) give
\(\gamma=0.2375330\) (\(\gamma/\gamma_{\min}=1.864700\), \(\beta=0.549406\)); all spins with \(g=2j+1\) give \(\gamma=0.2740669\)
(\(\gamma/\gamma_{\min}=2.151501\), \(\beta=0.604468\)).
\end{proposition}

\begin{proof}
The first case is \(2e^{-2\pi\gamma\sqrt3/2}=1\).  The other two were solved numerically and agree with~\cite{meissner} and~\cite{ghosh}
to the digits given; an independent recomputation agreed to \(10^{-12}\) (layer v4.9.0-r68).
\end{proof}

Under this reading \(\beta_{\rm sym}\) is not a fitted number but the ratio of two counting rules: the tachyon frame counts minimally
(\(j=\tfrac12\) only), the matter frame counts all spins with weight two.  One principle gives both rules:

\begin{proposition}[Uncertainty saturation]\label{prop:robertson}
In \(|j,m\rangle\) (\(\hbar=1\)), \(\Delta J_x\,\Delta J_y=\tfrac12\bigl(j(j+1)-m^2\bigr)\), and the Robertson bound
\(\Delta J_x\,\Delta J_y\ge\tfrac12|\langle J_z\rangle|=\tfrac12|m|\) holds with equality exactly for \(m=\pm j\).  At \(j=\tfrac12\) every state
saturates.
\end{proposition}

\begin{proof}
\(\langle J_x\rangle=\langle J_y\rangle=0\) and \(\langle J_x^2\rangle=\langle J_y^2\rangle=\tfrac12(j(j+1)-m^2)\).  The difference to the bound is
\(\tfrac12(j(j+1)-m^2-|m|)=\tfrac12(j-|m|)(j+|m|+1)\), which vanishes exactly for \(|m|=j\).  (Also formalized in Lean, module
\texttt{QSFormal/Zaehlung}.)
\end{proof}

So ``count only states that saturate the uncertainty bound'' reproduces both rules: the tachyon frame takes the absolute minimum, the
matter frame the two saturating states \(m=\pm j\) at every spin, which are exactly the states counted in~\cite{meissner}.  The Ghosh--Mitra
counting~\cite{ghosh} also counts non-saturating states.  A kinematic remark supports the minimal counting of the tachyon frame: a
tachyon that loses energy accelerates towards a single floor \(E=0\), and the floor of the mirror frame at \(-\eta\) reproduces
\(\gamma_L=\cosh2\eta\) as the distance between the frames (layer v4.9.0-r63).

\subsection{The limit of the counting reading}

The reading has a clear weakness.  The SU(2)-invariant treatment of the horizon~\cite{enp} is regarded as better founded than the U(1)
treatment, and it leads to \(\gamma\approx0.274\), the value with \(g=2j+1\).  With an SU(2) pair of countings one obtains
\(\beta=0.604\).  The value \(\beta\approx0.549\) used throughout therefore rests on the counting that the later literature considers the
less well founded (layer v4.9.0-r68).  The uncertainty reading is a reinterpretation of the U(1) counting, exact as mathematics, but not
part of its derivation.  And the check was not blind: the target window was known when it was made.

\subsection{An independent window}

A mixing model of the Companion, in which the frame rapidity enters the lepton mixing angles (Section~\ref{sec:readouts}), fixes
\(\eta\in[0.611,0.647]\), i.e.\ \(\beta_{\rm sym}\in[0.545,0.570]\) and \(\gamma_L\in[1.844,1.961]\), without using the Barbero--Immirzi values.
The input value lies inside; the Hopf value (\(\gamma_L=2\)) and the Tribonacci value (\(1.839\)) lie just outside, and so does the SU(2)
value \(\beta=0.604\).  This consistency is not independent evidence either: the reading of the model was chosen after the windows were
known (layer v4.9.0-r61).

\medskip\noindent\emph{Status.}  \(\beta_{\rm sym}\) is an input.  Propositions~\ref{prop:countings} and~\ref{prop:robertson} are exact.  That the
frames count by uncertainty saturation is a conjecture, and it selects the U(1) counting, which the literature does not favour.  OP1 reads:
derive from the frames why the tachyon frame counts only \(j=\tfrac12\) and the matter frame all spins with weight two --- or find a
different origin of \(\beta_{\rm sym}\).
